\documentclass[11pt]{article}
\usepackage{mathblatt}
% gesetzt von werkzeuge/setzer.py (Sorte selbst) aus dem Lernweg FS9
% und den Bankzeilen; Muster bau/proben/2026-10-09/hypotenuse-muster.
\usepackage{enumitem}
\usepackage{adjustbox,varwidth}
\newgeometry{a4paper,margin=18mm,top=15mm,bottom=20mm,footskip=9mm}
\pagestyle{fancy}\fancyhf{}\renewcommand{\headrulewidth}{0pt}
\fancyfoot[L]{\footnotesize\color{mbgrau}FS9 \textperiodcentered{} Lösungen \textperiodcentered{} Brüche vergleichen und ordnen}
\fancyfoot[R]{\footnotesize\color{mbgrau}\thepage}
\setlength{\parskip}{3pt}
\renewcommand{\kreuz}[1]{\mbox{$\square$\ #1}\quad}
\newcommand{\kreuzl}[1]{\par\hangindent1.4em\hangafter1\noindent$\square$\ #1}
\newcommand{\aufg}[3]{\par\noindent\makebox[0pt][r]{#3}\begin{minipage}[t]{\linewidth}#1\end{minipage}\par}
\newcommand{\aufgg}[4]{\par\noindent\makebox[0pt][r]{#4}\begin{minipage}[t]{0.6\linewidth}\vspace{0pt}#1\end{minipage}\hfill
  \begin{minipage}[t]{0.37\linewidth}\vspace{0pt}\centering #2\end{minipage}\par}
\newcommand{\nr}[1]{\makebox[7mm][l]{\textbf{#1}}}
\newcommand{\lsg}[3]{\par\noindent\begin{minipage}[t]{9mm}\textbf{#1}\end{minipage}%
  \begin{minipage}[t]{52mm}\raggedright\bfseries\boldmath #2\end{minipage}\hspace{3mm}%
  \begin{minipage}[t]{\dimexpr\linewidth-64mm\relax}\raggedright\small #3\end{minipage}\par\vspace{4pt}}
\newlength{\kr}
\begin{document}

\noindent{\Large\bfseries Lösungen}\hspace{1em}{\color{mbgrau}Brüche vergleichen und ordnen}\par\smallskip
\noindent{\small Links das Ergebnis, rechts der Weg in kurzen Schritten. Stimmt dein Ergebnis nicht, such den ersten Schritt, der anders ist.}\par
\par\Needspace{25mm}\vspace{6pt}\noindent\colorbox{black!12}{\makebox[7mm]{\bfseries\strut A}}\hspace{2mm}{\bfseries Gleicher Nenner oder gleicher Zähler}\par\vspace{4pt}
\lsg{1.}{a) $<$ \quad b) $>$ \quad c) $<$}{a) $\frac{4}{9} < \frac{7}{9}$; b) $\frac{11}{12} > \frac{5}{12}$; c) $\frac{6}{10} < \frac{9}{10}$}
\lsg{2.}{a) $>$ \quad b) $<$ \quad c) $<$}{a) $\frac{2}{3} > \frac{2}{7}$; b) $\frac{5}{12} < \frac{5}{6}$; c) $\frac{7}{10} < \frac{7}{8}$}
\lsg{3.}{$\frac{1}{10} < \frac{1}{6} < \frac{1}{4} < \frac{1}{3}$}{(gleicher Zähler: je größer der Nenner, desto kleiner)}
\lsg{4.}{Jan am weitesten, Ole am wenigsten weit}{gleicher Zähler, $\frac{3}{4} > \frac{3}{5} > \frac{3}{8}$ (Viertel sind die größten Teile, Achtel die kleinsten).}
\par\Needspace{25mm}\vspace{6pt}\noindent\colorbox{black!12}{\makebox[7mm]{\bfseries\strut B}}\hspace{2mm}{\bfseries Gleichnamig machen, dann vergleichen}\par\vspace{4pt}
\lsg{1.}{a) $<$ \quad b) $>$ \quad c) $>$}{a) $\frac{2}{3} = \frac{4}{6} < \frac{5}{6}$; b) $\frac{3}{4} = \frac{6}{8} > \frac{5}{8}$; c) $\frac{3}{5} = \frac{6}{10}$, also $\frac{7}{10} > \frac{3}{5}$}
\lsg{2.}{a) $<$ \quad b) $>$ \quad c) $>$}{a) $\frac{2}{3} = \frac{8}{12} < \frac{9}{12} = \frac{3}{4}$; b) $\frac{4}{5} = \frac{16}{20} > \frac{15}{20} = \frac{3}{4}$; c) $\frac{5}{6} = \frac{15}{18} > \frac{14}{18} = \frac{7}{9}$}
\lsg{3.}{a) $<$ \quad b) $>$ \quad c) $<$}{a) keins: $\frac{3}{4} = \frac{15}{20} < \frac{18}{20} = \frac{9}{10}$; b) gleicher Zähler: $\frac{5}{7} > \frac{5}{9}$; c) keins: $\frac{2}{5} = \frac{14}{35} < \frac{15}{35} = \frac{3}{7}$}
\lsg{4.}{Nein, es fehlt 1 Stimme.}{Nein; zugestimmt haben $\frac{17}{24}$, nötig sind $\frac{3}{4} = \frac{18}{24}$. Es fehlt 1 Stimme.}
\par\Needspace{25mm}\vspace{6pt}\noindent\colorbox{black!12}{\makebox[7mm]{\bfseries\strut C}}\hspace{2mm}{\bfseries Am Zahlenstrahl: ablesen, eintragen, ordnen}\par\vspace{4pt}
\lsg{1.}{$A = \frac{2}{5}$, \ $B = \frac{4}{5}$, \ $C = \frac{7}{5}$}{Ein Schritt ist $\frac{1}{5}$: $A = \frac{2}{5}$, $B = \frac{4}{5}$, $C = \frac{7}{5}$}
\lsg{2.}{$\frac{4}{12}$, \ $\frac{9}{12}$, \ $\frac{10}{12}$}{$\frac{1}{3} = \frac{4}{12}$ (4. Strich), $\frac{3}{4} = \frac{9}{12}$ (9. Strich), $\frac{5}{6} = \frac{10}{12}$ (10. Strich)}
\lsg{3.}{$\frac{3}{7}, \frac{2}{5}$ \,|\, $\frac{5}{9}, \frac{6}{11}$ \,|\, $\frac{8}{7}, \frac{11}{10}$}{kleiner als $\frac{1}{2}$: $\frac{3}{7}$, $\frac{2}{5}$; zwischen $\frac{1}{2}$ und 1: $\frac{5}{9}$, $\frac{6}{11}$; größer als 1: $\frac{8}{7}$, $\frac{11}{10}$}
\lsg{4.}{$\frac{3}{8} < \frac{2}{5} < \frac{5}{6} < \frac{7}{4}$}{zwei liegen unter $\frac{1}{2}$: $\frac{3}{8} = \frac{15}{40} < \frac{16}{40} = \frac{2}{5}$}
\lsg{5.}{$C < A < B < D$; $D$ am vollsten}{$D$ am vollsten; Reihenfolge $C < A < B < D$: $C = \frac{36}{80} = \frac{9}{20}$, $A = \frac{30}{50} = \frac{3}{5} = \frac{24}{40}$, $B = \frac{45}{72} = \frac{5}{8} = \frac{25}{40}$, $D = \frac{70}{80} = \frac{7}{8}$}
\par\Needspace{25mm}\vspace{6pt}\noindent\colorbox{black!12}{\makebox[7mm]{\bfseries\strut D}}\hspace{2mm}{\bfseries Einen Bruch dazwischen finden}\par\vspace{4pt}
\lsg{1.}{a) z.\,B. $\frac{3}{9}, \frac{4}{9}$ \quad b) z.\,B. $\frac{6}{11}, \frac{7}{11}$}{a) z.\,B. $\frac{3}{9}$ und $\frac{4}{9}$ (auch $\frac{5}{9}$); b) z.\,B. $\frac{6}{11}$ und $\frac{7}{11}$ (auch $\frac{8}{11}$)}
\lsg{2.}{a) $\frac{7}{16}$ \quad b) $\frac{3}{10}$}{a) $\frac{7}{16}$ (mit 2 erweitert: $\frac{6}{16}$ und $\frac{8}{16}$); b) $\frac{3}{10}$ (aus $\frac{2}{10}$ und $\frac{4}{10}$)}
\lsg{3.}{a) z.\,B. $\frac{1}{2}$ \quad b) $\frac{9}{20}$}{a) $\frac{1}{3} = \frac{4}{12}$, $\frac{3}{4} = \frac{9}{12}$, z.\,B. $\frac{6}{12} = \frac{1}{2}$; b) $\frac{4}{10}$ und $\frac{5}{10}$, mit 2: $\frac{8}{20}$ und $\frac{10}{20}$, dazwischen $\frac{9}{20}$}
\lsg{4.}{$\frac{5}{12}$}{$\frac{1}{4} = \frac{3}{12}$ und $\frac{2}{3} = \frac{8}{12}$}
\lsg{5.}{z.\,B. $\frac{22}{100}$\,m $=$ 22\,cm}{z.\,B. $\frac{22}{100}$\,m = 22\,cm; mit Nenner 100: $\frac{1}{5} = \frac{20}{100}$, $\frac{1}{4} = \frac{25}{100}$ (jede Länge zwischen 20\,cm und 25\,cm passt)}
\par\Needspace{25mm}\vspace{6pt}\noindent\colorbox{black!12}{\makebox[7mm]{\bfseries\strut T}}\hspace{2mm}{\bfseries Probetest}\par\vspace{4pt}
\lsg{1.}{a) $>$ \quad b) $<$ \quad c) $>$}{a) $\frac{5}{9} > \frac{4}{9}$; b) $\frac{2}{7} < \frac{2}{5}$; c) $\frac{3}{4} = \frac{15}{20} > \frac{14}{20} = \frac{7}{10}$}
\lsg{2.}{$A = \frac{3}{4}$, \ $B = \frac{5}{4}$, \ $\frac{3}{2} = \frac{6}{4}$}{Ein Schritt ist $\frac{1}{4}$: $A = \frac{3}{4}$, $B = \frac{5}{4}$; $\frac{3}{2} = \frac{6}{4}$ (6. Strich)}
\lsg{3.}{$\frac{1}{3} < \frac{1}{2} < \frac{4}{7} < \frac{5}{4}$}{}
\lsg{4.}{z.\,B. $\frac{2}{3}$}{$\frac{3}{5} = \frac{18}{30}$, $\frac{5}{6} = \frac{25}{30}$, $\frac{2}{3} = \frac{20}{30}$ liegt dazwischen}
\lsg{5.}{$\frac{7}{12} > \frac{1}{2}$}{denn $\frac{1}{2} = \frac{6}{12}$}
\lsg{6.}{die 6b}{Die 6b; $\frac{21}{28} = \frac{3}{4} = \frac{15}{20}$, $\frac{20}{25} = \frac{4}{5} = \frac{16}{20}$}
\end{document}
